% Section 6 -- Counterfactual policies (outline: paragraphs and what each must
% say; no prose yet). Written 29 October 2026.
%
% Length rule (agreed 29 October): the main text runs to about 40--45 pages
% with about 12 exhibits. Each result goes either to the main text, with full
% intuition, or to the online appendix, as a table and one paragraph with a
% one-sentence pointer from the text. Budget for this section: about 8 pages
% and 1--2 exhibits. \tier{MAIN} and \tier{ONLINE APP.} mark the tier of each
% item.
%
% Numbers are from the draft exhibits of the replication package: the policy
% network ("island"), seed 1, n = 65,536, 30 September batch. None of these
% tables is in the paper yet.
%   levers     output/tables/tab-policy-levers.csv   (m5 factorial, m7 doses)
%   variants   output/tables/tab-policy-variants.csv (m4: mandate, energy
%              prices, bill gap, mortality, no A-C for sale; island vs R0)
%   seeds      output/tables/tab-policy-seeds.csv    (m5 and m6)
%   doses      output/figures/fig-policy-dose.csv    (m7)
% The Shapley split by age and the flows by route and age were computed on 29
% October from the m5 and m6 cells in a scratch script. They are not yet in
% any package script.
% Decided 29 October: the section reports only changes from the policy
% network's own baseline, never its levels, and the network is not retrained.
% Every level in the paper comes from the estimated model (Section 5). So no
% ratio of shares and no statement that needs a level (such as the gradient
% reversing); proportional statements use percent changes within the network.
\providecommand{\tier}[1]{\textsf{\small [#1]}~}

% The opening and 6.1 were drafted as prose in 06-counterfactual-policies.tex
% on 5 November 2026; their outline is kept here but not typeset.
\iffalse
\outpara{Opening: the question and the answer (before the first subsection)}
\begin{itemize}
  \item \tier{MAIN} The question: which policies raise efficiency among older
    owners? Section~\ref{sec:results} implies that the route a policy acts on
    decides whom it reaches. Older owners rarely buy, so a policy that acts
    at a purchase reaches mostly the young. A policy that lowers the cost of
    renovating in place acts on the threshold gap, which every owner faces.
  \item \tier{MAIN} The answer up front. Halving the variable cost of
    renovation raises the A--C share of owners at 65--74 by 8.1 points and at
    25--34 by 2.5, narrowing the gap between them by 5.6 points, more than
    half the 9.6-point gap at the estimates (65--75; to be restated in the
    same band after the re-simulation). A point-of-sale standard
    raises the share the most, by 13.2 points, but widens the gap by 6.4.
  \item \tier{MAIN} Road map: 6.1 to 6.4.
  \item \open{Title: ``Counterfactual Policies'' or one that states the
    question, e.g.\ ``Which Policies Reach Older Owners''.}
\end{itemize}

\subsection{How the counterfactuals are computed}

\outpara{One paragraph in the main text; details in the online appendix}
\begin{itemize}
  \item \tier{MAIN} One network holds the estimated parameters fixed and
    takes the policy parameters as inputs: $\xi_f$, $\xi$, $\zeta_b$,
    $\zeta_s$, $\lambda_L$, $\lambda_D$, each trained over a box around its
    baseline. Each policy is then a simulation in which households
    re-optimize. The point-of-sale standard and the energy-price variants
    change the model, not a lever, so each is trained from the same network.
  \item \tier{MAIN} Every policy uses the same draws as the baseline (common
    random numbers). Effects are differences from that baseline, in
    percentage points.
  \item \tier{MAIN} Scope: (i) partial equilibrium: house prices, the price
    premium of efficient homes, and the supply of efficient homes for sale
    are held fixed, which matters most for a large standard; (ii) each
    policy is in place over the whole life cycle, from 23 on, and the
    transition after a reform is not computed; (iii) no welfare
    calculation.
  \item \tier{MAIN, one sentence} The policy network is a second
    approximation of the same model, and its baseline lies below the
    estimates: 13.6 percent of all owners against 16.2. So we report only
    changes from its own baseline, and every level in the paper comes from
    the estimated model. Changes carry over: the standard imposed directly
    widens the gap by 4.1 points on each network.
  \item \tier{ONLINE APP.} Where the two baselines differ: both solve the
    same problem at the same parameters, and the two policy-network seeds
    agree; they differ mainly in whether a renter who buys picks an A--C
    home, a close call that small approximation errors tip. A table of the
    two baselines (shares, gap, ownership, moves), the cross-check of the
    standard, the per-age accuracy audits, and the training box.
    Appendix~\ref{app:computation} covers the network.
\end{itemize}
\fi

% Drafted as prose in 06-counterfactual-policies.tex on 5 November 2026
% (heading, text and Table tab:policy-renovation moved there). The outline
% below is kept for reference; its ``why the old gain most'' bullet is
% superseded: 72 against 22 D--G owners per 100 households reflects
% ownership by age, not the share among owners (82 and 88 percent of
% owners).
\iffalse
\subsection{Lowering the cost of renovating in place}

\outpara{The two costs and what they act on}
\begin{itemize}
  \item \tier{MAIN} The variable cost $\xi$ is halved from 0.08 to 0.04, and
    the fixed cost $\xi_f$ from 0.159 to 0.079. Either lowers the owner's
    threshold $c_j^*$ and narrows the threshold gap
    (Section~\ref{sec:results-window}).
\end{itemize}

\outpara{The variable cost halved (main table)}
\begin{itemize}
  \item \tier{MAIN} All owners $+6.3$; 25--34 $+2.5$; 65--74 $+8.1$. The gap
    narrows by 5.6 points.
  \item \tier{MAIN} Why the old gain most: renovation in place at 65--74
    rises by 2.0 per 100 households a year, from 0.9, and at 25--34 by 0.5,
    from 0.2. The subsidy acts on owners of D--G homes, and they are
    concentrated among the old: 72 per 100 households at 65--74 against 22
    at 25--34. Buying an efficient home hardly changes, by $-0.1$ at both
    ages.
  \item \tier{MAIN} Outlay: 5.7 ($\times10^3$, model money per
    household-year, inframarginal payments included), or 0.91 per point of
    the all-owner share.
\end{itemize}

\outpara{The fixed cost halved (main table)}
\begin{itemize}
  \item \tier{MAIN} All owners $+1.7$; 25--34 $+0.9$; 65--74 $+1.5$; gap
    $-0.6$. Same direction as the variable cost, smaller. Outlay
    0.73 per point.
  \item \open{Why the fixed cost does so much less than the variable cost
    at a similar outlay per point. Check against the size of each cost
    relative to a typical renovation before writing an intuition.}
\end{itemize}

\outpara{Dose (online appendix, one sentence in the text)}
\begin{itemize}
  \item \tier{ONLINE APP.} Figure fig:policy-dose. Effects grow more than in
    proportion to the cut. A cut of the variable cost to 25 percent raises
    the all-owner share by 13.6 points and narrows the gap by 10.6. Whether
    the gradient would reverse is a statement about levels, so it is not
    made.
\end{itemize}
\fi

% 6.2 drafted in 06-counterfactual-policies.tex (5 November 2026).
\iffalse
\subsection{Acting at a purchase}\label{sec:policy-purchase}

\outpara{A point-of-sale standard (main table)}
\begin{itemize}
  \item \tier{MAIN} Every purchase of a D--G home is renovated at purchase.
    All owners $+13.2$, the largest effect of any policy; 25--34 $+16.3$;
    65--74 $+9.9$. The gap widens by 6.4 points. Ownership falls by 0.3
    points, and owner moves fall by 2.1 per 100
    owners.
  \item \tier{MAIN} Why: the standard acts on the window, and the young
    pass through it most often. The old gain largely because they bought
    under the standard earlier in life. This comes from comparing whole life
    cycles (6.1), so a reform introduced today would reach today's old
    owners less.
  \item \open{Outlay: the standard costs the government nothing, and the
    cost falls on buyers. Leave the outlay column blank or report the
    buyers' renovation spending?}
\end{itemize}

\outpara{Buying and selling costs halved (main table, level only)}
\begin{itemize}
  \item \tier{MAIN} All owners $+4.9$. Owner moves rise by 10.0 per 100
    owners, from 7.6. More owners pass through the window: buying an
    efficient home rises by 2.6 per 100 households a year at 25--34 and by
    3.0 at 65--74, while renovation in place falls slightly.
  \item \tier{MAIN} Outlay 5.46 per point, six times that of the variable
    cost.
  \item \tier{MAIN, one sentence; ONLINE APP.} The split by age is not
    robust: the gap changes by $+3.8$ on seed 1 and $+0.5$ on seed 2, so
    only the level is reported (Table tab:policy-seeds).
\end{itemize}
\fi

% 6.3 drafted in 06-counterfactual-policies.tex (5 November 2026; energy
% prices). The rest of this block (the online-appendix variants, the
% combined cost levers) is kept for the online appendix.
\iffalse
\subsection{Changing the payoff}\label{sec:policy-payoff}

\outpara{What shifts the level without tilting it (one paragraph)}
\begin{itemize}
  \item \tier{MAIN} Energy prices doubled: all owners $+5.2$; 25--34
    $+7.1$; 65--74 $+4.3$. Ownership falls by 2.9 points. In percent of the
    policy network's own baseline, the share rises by about 40 percent at
    25--34 and 37 percent at 65--74. Like comfort in
    Figure~\ref{fig:mechanisms-payoff}, a higher payoff per year raises the
    share at every age in about the same proportion.
  \item \tier{ONLINE APP., one sentence each} No energy-bill gap between
    A--C and D--G homes: $-3.5$. Credit limits: debt-to-income
    0.33 to 0.365, 0.0; loan-to-value 0.90 to 0.85, $-0.2$. Mortality halved
    and no A--C homes for sale on the policy network, as a check on
    Section~\ref{sec:results}.
\end{itemize}

\outpara{Combining the cost levers (online appendix, one sentence in the
text)}
\begin{itemize}
  \item \tier{ONLINE APP.} All three halved: $+12.2$ (seed 2: $+12.7$).
    Shapley shares: variable cost 57 percent, fixed cost 26, buying and
    selling costs 17 (seed 2: 54, 25, 21).
  \item \tier{ONLINE APP.} The levers interact. The fixed cost's share, 3.2
    points, is twice its effect alone, 1.7. The share of buying and selling
    costs, 2.1, is less than half its effect alone, 4.9. Cheaper renovation
    makes the fixed cost bite more and a cheaper move matter less.
  \item \tier{ONLINE APP.} Split of the gap: the variable cost $-5.5$ and
    $-5.1$ on the two seeds, the fixed cost $-1.5$ and $-1.4$; buying and
    selling costs $+4.8$ and $+1.6$, not robust.
\end{itemize}
\fi

% Decided 5 November 2026: the section describes the implications of each
% policy and does not rank them. So no outlay per point (report the total
% outlay), and the close below must be rewritten without ``lowest outlay
% per point'' or other ranking language.
\outpara{Section close (one paragraph)}
\begin{itemize}
  \item \tier{MAIN} Whom a policy reaches depends on the route it acts on.
    Lowering the cost of renovating in place reaches older owners and
    flattens the gradient, at the lowest outlay per point. Acting at a
    purchase raises the level the most but widens the gradient. Raising the
    payoff per year raises the share in about the same proportion at every
    age.
  \item \tier{MAIN} Scope, restated briefly: partial equilibrium, whole life
    cycles, no welfare.
\end{itemize}

\outpara{Exhibits and open decisions (planning note, not text)}
\begin{itemize}
  \item \tier{MAIN} One table, tab:policy, of changes only (no level row;
    the notes point to the levels at the estimates in Section 5), with
    numbered columns built in the R script: (1) all owners, (2) 25--34,
    (3) 65--74, (4) gap, (5) owner moves, (6) outlay per point. Rows:
    \begin{itemize}
      \item renovating in place: the variable cost halved, the fixed cost
        halved;
      \item at a purchase: the point-of-sale standard (column 6 a dash: a
        regulation, no public outlay), buying and selling costs halved
        (columns 1, 5 and 6 only);
      \item the payoff: energy prices doubled.
    \end{itemize}
    It merges \texttt{tab\_policy\_levers.R} and
    \texttt{tab\_policy\_variants.R}.
  \item \open{An optional main figure, which would be new: the effect at
    65--74 against the effect at 25--34 for each policy, with the
    45-degree line, so that which policies tilt the profile is visible at
    a glance.}
  \item \tier{ONLINE APP.} Table tab:policy-seeds, extended with the Shapley
    split by age; credit limits, the bill gap, mortality halved and no A--C
    homes for sale (tab:policy-variants-other, extended); the two baselines
    and the cross-check of the standard; Figure fig:policy-dose.
  \item \open{Experiments not yet run, each needing a training on a pod:
    a grant for renovating in place only; a larger price premium through
    disclosure ($\gamma_p$){\draft, in training since 5 November 2026}.
    None is needed for the argument above.}
  \item \open{Welfare by age (from the value networks, new code) or energy
    saved (from the simulations): a JPE referee is likely to ask for one
    of them.}
\end{itemize}
